\begin{abstract}
本文提出一种灵活的新技术，用于方便地完成相机标定。该方法仅要求相机在少数（至少两个）不同方向下观测一个平面模板。相机或平面模板均可自由移动，且不需要知道其运动。方法对镜头的径向畸变进行建模。所提出的过程由一个闭式解开始，随后依据最大似然准则进行非线性细化。我们使用计算机仿真数据和真实数据对所提出的技术进行了测试，并获得了非常好的结果。与使用昂贵设备（例如两个或三个相互正交的平面）的传统技术相比，所提出的技术易于使用且具有灵活性。它使三维计算机视觉从实验室环境向现实应用迈进一步。相应的软件可从作者的网页获得。
\end{abstract}
\keywords{相机标定；内参数；镜头畸变；基于平面的灵活标定；运动分析；模型获取}
